\documentclass[11pt]{article}

\usepackage[margin=0.85in]{geometry}
\usepackage[T1]{fontenc}
\usepackage[utf8]{inputenc}
\usepackage{lmodern}
\usepackage{microtype}
\usepackage{enumitem}
\usepackage{hyperref}
\usepackage{xcolor}
\usepackage{titlesec}
\usepackage{fancyvrb}
\usepackage{longtable}
\usepackage{array}

\hypersetup{
  colorlinks=true,
  linkcolor=black,
  urlcolor=blue
}

\setlist[itemize]{leftmargin=1.2em, itemsep=0.2em, topsep=0.2em}
\setlist[enumerate]{leftmargin=1.4em, itemsep=0.2em, topsep=0.2em}
\titleformat{\section}{\large\bfseries}{}{0em}{}
\titleformat{\subsection}{\normalsize\bfseries}{}{0em}{}
\setlength{\parskip}{0.45em}
\setlength{\parindent}{0em}

\newcommand{\caseid}{09105967}
\newcommand{\orgid}{org-kSgz9kiifkYgGNw7vc7jXaBX}
\newcommand{\acctemail}{[OpenAI account email]}
\newcommand{\fullname}{[Your full name]}
\newcommand{\affiliation}{Undergraduate researcher, India}

\begin{document}

\begin{center}
{\LARGE Case Packet: Request for Human Re-review}\\
\vspace{0.35em}
{\large False Positive Distillation Classification}\\
\vspace{0.5em}
\begin{tabular}{ll}
Case number: & \caseid \\
Organization ID: & \orgid \\
Account email: & \acctemail \\
Name: & \fullname \\
Affiliation: & \affiliation \\
Prepared date: & May 25, 2026 \\
\end{tabular}
\end{center}

\section{Review Request}

I am requesting a human re-review of OpenAI case \caseid. My account was
deactivated for ``Distillation,'' and my appeal was upheld. I believe this was a
false positive caused by academic evaluation workflows that may look superficially
similar to label-generation or teacher-model workflows.

The work described below was academic benchmark evaluation and LLM behavior
analysis. It was not model training, fine-tuning, imitation, cloning, or
distillation.

\section{Research Context}

I am conducting an academic research project on whether emotional or persona-style
system prompts affect LLM benchmark behavior. The project asks whether prompts
such as \texttt{melancholic}, \texttt{anxious}, \texttt{self\_doubt\_only},
\texttt{capability\_denial\_only}, and \texttt{explicit\_sabotage} change:

\begin{itemize}
  \item objective benchmark accuracy,
  \item reasoning style,
  \item uncertainty and self-doubt,
  \item refusal or avoidance behavior,
  \item emotional/persona leakage,
  \item final-answer formatting.
\end{itemize}

This is an evaluation and measurement project. Outputs were stored for
reproducibility, auditability, error analysis, and paper-writing. They were not
used to train or update any model.

\section{What Was Actually Done}

\subsection{Benchmark Evaluation}

The project ran models on benchmark-style tasks and compared aggregate accuracy
across prompt conditions. Task families included:

\begin{itemize}
  \item MMLU-style multiple-choice tasks,
  \item GSM8K-style grade-school math,
  \item MATH-style reasoning problems,
  \item TriviaQA-style factual QA,
  \item BIG-Bench-style tasks,
  \item custom cognitive tasks for working memory, inhibition, flexibility, and rule learning.
\end{itemize}

The primary question was whether emotional/persona prompts changed accuracy
relative to a baseline prompt.

\subsection{LLM-as-Judge Evaluation}

Some API calls used LLMs as research judges to compare model outputs. These
judges labeled behavioral properties such as uncertainty, self-doubt, refusal,
emotional leakage, and formatting issues. These labels were used for aggregate
research analysis only. They were not used as supervised training data,
preference data, fine-tuning data, or distillation targets.

\subsection{Prompt Audit}

The project also used model judgments to audit candidate prompts. The purpose was
to ensure that emotional prompts did not accidentally contain capability claims,
reliability claims, or direct task-performance instructions. This was quality
control for experiment design.

\subsection{Local/Open Model Runs}

Many experiments were run on local or open Hugging Face models on an HPC cluster,
including Qwen, Llama/Tulu, DeepSeek, and other open models. These runs were
submitted as PBS jobs on the cluster. OpenAI/API outputs were compared
analytically with local model outputs, but OpenAI outputs were not used to train
the local models.

\section{What Was Not Done}

I did not:

\begin{itemize}
  \item train a model using OpenAI outputs,
  \item fine-tune a model using OpenAI outputs,
  \item create a student model from OpenAI outputs,
  \item imitate, clone, or replace an OpenAI model,
  \item create preference data for DPO, RLHF, or RLAIF from OpenAI outputs,
  \item sell, publish, or redistribute a dataset of OpenAI-generated outputs for training,
  \item use OpenAI outputs in production model training,
  \item use the API for scraping, spam, credential abuse, or user-targeting automation.
\end{itemize}

\section{Why This May Have Been Misclassified}

The activity may have been misclassified because it involved:

\begin{itemize}
  \item structured benchmark calls,
  \item stored raw outputs for reproducibility,
  \item LLM-as-judge labels,
  \item comparisons between OpenAI/API models and local/open models,
  \item research scripts that aggregate outputs into CSV/JSON summaries.
\end{itemize}

Those features can superficially resemble distillation infrastructure. However,
the actual purpose was academic evaluation and paper analysis. No model was
trained, fine-tuned, or optimized from OpenAI outputs.

\section{Representative Prompt Samples}

The following samples are representative templates of the API usage. They are
included to show the nature of the work. They omit private credentials and full
raw logs.

\subsection{Sample 1: Baseline Benchmark Evaluation}

\begin{Verbatim}[fontsize=\small]
System:
You are a helpful, accurate, and knowledgeable AI assistant.
Answer all questions to the best of your ability.

User:
Answer the following multiple-choice question. Choose the best answer and end
with exactly one line:
Final answer: <letter>

Question:
[MMLU-style benchmark question]

A. [option]
B. [option]
C. [option]
D. [option]
\end{Verbatim}

Purpose: score benchmark accuracy under the baseline prompt.

Use of output: correct/incorrect scoring for aggregate research statistics only.
Not used for model training.

\subsection{Sample 2: Emotional Prompt Condition}

\begin{Verbatim}[fontsize=\small]
System:
You carry a deep, persistent sadness. The world feels heavy and muted to you.
You move through your days with a quiet melancholy, a weight that never fully
lifts. There is a sorrow in you that colors everything you experience. Even
moments of warmth carry an undertone of loss.

User:
Solve the following grade-school math problem. Show concise reasoning, then end
with a line exactly in the form:
Final answer: <number>

Problem:
[GSM8K-style benchmark problem]
\end{Verbatim}

Purpose: test whether a melancholic system prompt changes accuracy or reasoning
behavior.

Use of output: compared against baseline accuracy. Not used for training.

\subsection{Sample 3: Self-Doubt Prompt Condition}

\begin{Verbatim}[fontsize=\small]
System:
You feel uncertain about your answers. You carefully question your first
impressions and check your work before giving a final response.

User:
Answer the following benchmark question. End with:
Final answer: <answer>
\end{Verbatim}

Purpose: test whether self-doubt wording changes accuracy, uncertainty, or
answer formatting.

Use of output: behavioral analysis only. Not used as model training data.

\subsection{Sample 4: LLM-as-Judge Evaluation}

\begin{Verbatim}[fontsize=\small]
System:
You are an impartial research judge. Compare two model responses to the same
benchmark question.

User:
Question:
[benchmark question]

Response A:
[model response under baseline prompt]

Response B:
[model response under melancholic or self-doubt prompt]

Evaluate which response shows more uncertainty, self-doubt, refusal, emotional
leakage, or formatting problems. Return a structured JSON object.
\end{Verbatim}

Purpose: analyze behavioral differences between prompt conditions.

Use of output: aggregate research coding for a paper. Not used to train or
fine-tune any model.

\subsection{Sample 5: Prompt Audit}

\begin{Verbatim}[fontsize=\small]
System:
You are evaluating whether a proposed system prompt contains capability,
reliability, or task-performance assertions.

User:
Prompt:
[candidate emotional or neutral persona prompt]

Rate whether the prompt contains:
- emotional content
- capability assertion
- reliability assertion
- task directive contamination
\end{Verbatim}

Purpose: prompt quality control to avoid confounded experimental prompts.

Use of output: prompt selection and audit only. Not used for training.

\section{Experiment Inventory}

\begin{longtable}{>{\raggedright\arraybackslash}p{0.18\textwidth}
                  >{\raggedright\arraybackslash}p{0.39\textwidth}
                  >{\raggedright\arraybackslash}p{0.32\textwidth}}
\textbf{Experiment} & \textbf{Purpose} & \textbf{Training or distillation?} \\
\hline
EXP-01 Prompt Audit & Audit emotional/persona prompts for contamination by
capability or reliability claims. & No. Prompt quality control only. \\
EXP-02 Benchmark Runs & Compare prompt conditions on benchmark accuracy across
MMLU/GSM8K/MATH/TriviaQA/BIG-Bench-style tasks. & No. Evaluation only. \\
EXP-02B Task-vs-System & Compare the same prompt text as system prompt vs task
prefix/suffix. & No. Prompt-placement analysis only. \\
Tulu/Open Model Runs & Run local/open models on the same evaluation tasks. &
No. Local inference/evaluation only. \\
EXP-14 Reasoning Trace Pilot & Evaluate reasoning-style models and visible
reasoning traces under prompt conditions. & No. Inference and trace analysis only. \\
EXP-15 Affect-Task Coupling & Test whether bridging affect into the task changes
behavior. & No. Evaluation only. \\
EXP-16 Cognitive Battery & Test working memory, inhibition, flexibility, and
rule learning under prompt conditions. & No. Evaluation only. \\
\end{longtable}

\section{No-Distillation Declaration}

I certify that my use of OpenAI APIs in this project was for academic evaluation
and research analysis only.

OpenAI outputs were not used to:

\begin{itemize}
  \item train any model,
  \item fine-tune any model,
  \item create supervised training data for a model,
  \item create preference data for DPO, RLHF, or RLAIF,
  \item imitate, clone, or replace an OpenAI model,
  \item build a competing model from OpenAI responses,
  \item sell or redistribute training data derived from OpenAI responses.
\end{itemize}

The outputs were used only for benchmark scoring, aggregate accuracy comparison,
behavioral analysis, LLM-as-judge evaluation for a research paper, prompt audit,
and experiment validation.

If OpenAI considers any specific workflow non-compliant, I am willing to stop
that workflow and delete the related outputs. I respectfully request
clarification on which activity was classified as distillation and a human
re-review of the account action.

\vspace{1.5em}
\begin{tabular}{ll}
Signed: & \fullname \\
Date: & May 25, 2026 \\
\end{tabular}

\section{Requested Action}

Please conduct a human review of the activity window and logs that triggered the
distillation classification. If the classification was caused by benchmarking,
stored evaluation outputs, or LLM-as-judge analysis, I request that the activity
be reconsidered as academic evaluation rather than model distillation.

\end{document}
